\begin{center}\textbf{چکیده}\end{center}
\noindent
آموزش هم‌آورد از روش‌های مؤثر برای افزایش ستبری شبکه‌های عصبی در برابر آشفتگی‌های هدفمند ورودی است؛ بااین‌حال، ساخت نمونه‌های هم‌آورد با حملات چندگامی هزینه محاسباتی زیادی دارد. روش‌های سریع با تقریب تک‌گامی مسئله بیشینه‌سازی درونی این هزینه را کاهش می‌دهند، اما ممکن است به بیش‌برازش فاجعه‌بار دچار شوند: درستی در برابر حملات چندگامی به‌طور ناگهانی فرو می‌ریزد، درحالی‌که درستی در برابر حمله تک‌گامی همچنان زیاد است. این پایان‌نامه مسئله پیش‌بینی و پیشگیری کم‌هزینه از این پدیده را با بررسی هندسه تابع هزینه مطالعه می‌کند.

نخست، بیش‌برازش اپسیلون به‌عنوان حساسیت مدل به اندازه‌ها و جهت‌های محدود آشفتگی بررسی می‌شود. این دیدگاه توضیح می‌دهد چگونه آموزش روی نقاط حاصل از یک گام ثابت می‌تواند ناحیه‌های آسیب‌پذیر میان ورودی و نمونه هم‌آورد را پنهان کند. سپس شاخص هم‌راستایی آشفتگی، موسوم به \pertalign، از شباهت کسینوسی شیب ورودی پیش و پس از حمله تعریف می‌شود. تحلیل تیلور نشان می‌دهد افت این شاخص با مؤلفه عمود بر شیبِ اثر هسین در جهت آشفتگی ارتباط دارد. استفاده مجدد از شیب‌های موجود در آموزش، پایش این شاخص را بدون محاسبه صریح هسین و با سربار اندک ممکن می‌سازد.

بر پایه این تحلیل، روش تطبیقی مرتبه دوم \sora{} معرفی می‌شود که با تخمین خمیدگی جهتی، میانگین‌گیری نمایی و نمونه‌برداری مؤلفه‌ای اندازه گام، توزیع آشفتگی را در طول آموزش تنظیم می‌کند. نتایج منتشرشده پژوهش روی دادگان تصویری طبیعی و زیست‌پزشکی و معماری‌های گوناگون، بهبود بده‌بستان درستی پاک، ستبری و هزینه را نشان می‌دهند. برای نمونه، در ارزیابی \lr{AutoAttack} با بودجه $8/255$ روی \lr{PreActResNet-18}، درستی \sora{} در \lr{PathMNIST} برابر $32.91\%$ و در \lr{ImageNet-100} برابر $18.56\%$ گزارش شده است. این اعداد بازگویی نتایج مقاله مشترک منتشرشده در \lr{ICML 2026} هستند. همچنین، محدودیت تقریب محلی، تفاوت دامنه آشفتگی آموزش و ارزیابی، و نیاز به بررسی آماری هشدار زودهنگام تشریح می‌شود.

\bigskip
\noindent\textbf{کلیدواژه‌ها:}
آموزش هم‌آورد سریع، بیش‌برازش فاجعه‌بار، بیش‌برازش اپسیلون، هم‌راستایی آشفتگی، خمیدگی، ستبری شبکه عصبی.
